The finite-support hypothesis in Theorem~\ref{thm:finite-inverse-projectives}
can be removed if one describes the ordinary projective types by local
idempotents in the canonical action ideal. A unital ring is
\emph{semiperfect} if its identity is a finite sum of orthogonal
idempotents with local corner rings; the zero ring is allowed.

\begin{theorem}[amplitudes with semiperfect upper coefficients]
\label{thm:upper-semiperfect-amplitudes}
Let $T=(Q,F)$ be a unital inverse semiring such that $F(e)$ is
semiperfect for every $e\in U(Q)$. For any $Q$-semimodule $L$, put
$\mathcal A=\mathcal A_L$ and $\mathcal J=\mathcal J_L$. Consider pairs
$(l,f)$ with
\[
 0\ne f=f^2\in\mathcal J(l,l),\qquad
 f\mathcal A(l,l)f\text{ a local ring}.
\]
Declare $(l,f)\sim(k,g)$ if there are arrows $u:l\to k$ and
$v:k\to l$ in $\mathcal A$ satisfying
\[
 u=guf,\qquad v=fvg,\qquad vu=f,\qquad uv=g.
\]
Every finitely generated projective $\mathcal A$-module $G$ with
$\mathcal JG=G$ is a finite direct sum of
\[
 P_{l,f}=\mathcal A(l,-)f.
\]
Its multiplicities on the displayed equivalence classes are unique,
and every such finite sum occurs. Consequently finitely generated
projective $T$-semimodules are classified by a finitely generated
projective $Q$-support $L$ and these finite multiplicities, modulo
$\operatorname{Aut}_Q(L)$.

If every upper coefficient ring is local or zero, the permitted pairs
reduce to objects $l$ with $\mathcal A(l,l)\ne0$ and
$1_l\in\mathcal J(l,l)$, with $f=1_l$. The multiplicities are then
indexed by their isomorphism classes as objects of $\mathcal A$.
\end{theorem}

There is no finiteness assumption on $Q$, $L$ or the coefficient
groups, and no commutativity assumption. The empty sum is zero
amplitude over the given support. The theorem determines the
amplitude parameter in this class; it leaves the classification of
projective $Q$-supports and the evaluation of the displayed
idempotents and intertwiners as separate questions. In particular,
it does not assert that every action endomorphism ring is
semiperfect; Example~\ref{ex:upper-local-action-boundary} shows why.

